% ===== 1-D Dynamic Programming (12 problems) =====
\section{1-D Dynamic Programming \texorpdfstring{{\normalsize\color{muted}$\cdot$ 12 problems}}{- 12 problems}}
\markboth{1-D Dynamic Programming}{}

\begin{problem}{70}{Climbing Stairs}{Easy}
\summary{Ways to climb n steps taking 1 or 2 at a time.}
\keyidea{ways(n) = ways(n-1) + ways(n-2): Fibonacci with two variables.}
\complexity{O(n) time, O(1) space}
\begin{code}
class Solution:
    def climbStairs(self, n: int) -> int:
        a, b = 1, 1
        for _ in range(n - 1):
            a, b = b, a + b
        return b
\end{code}
\end{problem}

\begin{problem}{746}{Min Cost Climbing Stairs}{Easy}
\summary{Pay cost[i] to step from stair i (1 or 2 steps); start at 0 or 1; min cost to reach the top.}
\keyidea{dp[i] = cost[i] + min(dp[i+1], dp[i+2]); answer min(dp[0], dp[1]).}
\complexity{O(n) time, O(1) space}
\begin{code}
class Solution:
    def minCostClimbingStairs(self, cost: List[int]) -> int:
        a = b = 0                       # dp[i+1], dp[i+2]
        for i in range(len(cost) - 1, -1, -1):
            a, b = cost[i] + min(a, b), a
        return min(a, b)
\end{code}
\end{problem}

\begin{problem}{198}{House Robber}{Medium}
\summary{Max money from houses in a line without robbing two adjacent ones.}
\keyidea{best(i) = max(best(i-1), best(i-2) + nums[i]).}
\complexity{O(n) time, O(1) space}
\begin{code}
class Solution:
    def rob(self, nums: List[int]) -> int:
        prev2 = prev1 = 0
        for n in nums:
            prev2, prev1 = prev1, max(prev1, prev2 + n)
        return prev1
\end{code}
\end{problem}

\begin{problem}{213}{House Robber II}{Medium}
\summary{Same as House Robber but houses form a circle (first and last are adjacent).}
\keyidea{Run House Robber on nums[1:] and nums[:-1], take the max (handle n == 1).}
\complexity{O(n) time, O(1) space}
\begin{code}
class Solution:
    def rob(self, nums: List[int]) -> int:
        def line(arr):
            a = b = 0
            for n in arr:
                a, b = b, max(b, a + n)
            return b
        if len(nums) == 1:
            return nums[0]
        return max(line(nums[1:]), line(nums[:-1]))
\end{code}
\end{problem}

\begin{problem}{5}{Longest Palindromic Substring}{Medium}
\summary{Return the longest palindromic substring.}
\keyidea{Expand around every centre (odd: i,i and even: i,i+1).}
\complexity{O(n\textasciicircum{}2) time, O(1) space}
\begin{code}
class Solution:
    def longestPalindrome(self, s: str) -> str:
        res = ""
        for i in range(len(s)):
            for l, r in ((i, i), (i, i + 1)):
                while l >= 0 and r < len(s) and s[l] == s[r]:
                    l, r = l - 1, r + 1
                if r - l - 1 > len(res):
                    res = s[l + 1 : r]
        return res
\end{code}
\end{problem}

\begin{problem}{647}{Palindromic Substrings}{Medium}
\summary{Count all palindromic substrings.}
\keyidea{Expand around every centre; each successful expansion is one more palindrome.}
\complexity{O(n\textasciicircum{}2) time, O(1) space}
\begin{code}
class Solution:
    def countSubstrings(self, s: str) -> int:
        count = 0
        for i in range(len(s)):
            for l, r in ((i, i), (i, i + 1)):
                while l >= 0 and r < len(s) and s[l] == s[r]:
                    count += 1
                    l, r = l - 1, r + 1
        return count
\end{code}
\end{problem}

\begin{problem}{91}{Decode Ways}{Medium}
\summary{Digits map 1-\textgreater{}A ... 26-\textgreater{}Z. How many ways can the digit string be decoded?}
\keyidea{dp[i] = dp[i+1] (if s[i] != '0') + dp[i+2] (if s[i:i+2] in 10..26).}
\complexity{O(n) time, O(1) space}
\begin{code}
class Solution:
    def numDecodings(self, s: str) -> int:
        one, two = 1, 0                   # dp[i+1], dp[i+2]
        for i in range(len(s) - 1, -1, -1):
            cur = 0
            if s[i] != "0":
                cur = one
                if i + 1 < len(s) and 10 <= int(s[i : i + 2]) <= 26:
                    cur += two
            one, two = cur, one
        return one
\end{code}
\end{problem}

\begin{problem}{322}{Coin Change}{Medium}
\summary{Fewest coins (unlimited supply) to make amount; -1 if impossible.}
\keyidea{dp[a] = 1 + min(dp[a - c]) over coins c (unbounded knapsack).}
\complexity{O(amount * coins) time, O(amount) space}
\begin{code}
class Solution:
    def coinChange(self, coins: List[int], amount: int) -> int:
        dp = [0] + [amount + 1] * amount
        for a in range(1, amount + 1):
            for c in coins:
                if c <= a:
                    dp[a] = min(dp[a], 1 + dp[a - c])
        return dp[amount] if dp[amount] <= amount else -1
\end{code}
\end{problem}

\begin{problem}{152}{Maximum Product Subarray}{Medium}
\summary{Largest product of a contiguous subarray.}
\keyidea{Track both current max and min product (a negative can flip min into max).}
\complexity{O(n) time, O(1) space}
\begin{code}
class Solution:
    def maxProduct(self, nums: List[int]) -> int:
        res = max(nums)
        cur_max = cur_min = 1
        for n in nums:
            cand = (n, n * cur_max, n * cur_min)
            cur_max, cur_min = max(cand), min(cand)
            res = max(res, cur_max)
        return res
\end{code}
\end{problem}

\begin{problem}{139}{Word Break}{Medium}
\summary{Can s be split into a sequence of dictionary words (reuse allowed)?}
\keyidea{dp[i] = True if some word w matches s[i:] start and dp[i + len(w)] is True.}
\complexity{O(n * m * L) time, O(n) space}
\begin{code}
class Solution:
    def wordBreak(self, s: str, wordDict: List[str]) -> bool:
        dp = [False] * (len(s) + 1)
        dp[len(s)] = True
        for i in range(len(s) - 1, -1, -1):
            for w in wordDict:
                if s.startswith(w, i) and dp[i + len(w)]:
                    dp[i] = True
                    break
        return dp[0]
\end{code}
\end{problem}

\begin{problem}{300}{Longest Increasing Subsequence}{Medium}
\summary{Length of the longest strictly increasing subsequence.}
\keyidea{Patience sorting: tails[k] = smallest tail of an increasing subsequence of length k+1; binary search to place each number.}
\complexity{O(n log n) time, O(n) space}
\begin{code}
class Solution:
    def lengthOfLIS(self, nums: List[int]) -> int:
        tails = []
        for n in nums:
            i = bisect.bisect_left(tails, n)
            if i == len(tails):
                tails.append(n)
            else:
                tails[i] = n
        return len(tails)
\end{code}
\end{problem}

\begin{problem}{416}{Partition Equal Subset Sum}{Medium}
\summary{Can the array be split into two subsets with equal sums?}
\keyidea{Need a subset summing to total/2: keep a set (or bitset) of reachable sums.}
\complexity{O(n * sum) time, O(sum) space}
\begin{code}
class Solution:
    def canPartition(self, nums: List[int]) -> bool:
        total = sum(nums)
        if total % 2:
            return False
        target = total // 2
        reach = {0}
        for n in nums:
            reach |= {r + n for r in reach if r + n <= target}
            if target in reach:
                return True
        return target in reach
\end{code}
\end{problem}
